\documentclass[11pt]{article}
\usepackage[T1]{fontenc}
\usepackage{lmodern}
\usepackage{amsmath,amssymb,amsthm,mathtools}
\usepackage[a4paper,margin=1in]{geometry}
\usepackage[hidelinks]{hyperref}
\usepackage{microtype}

\newtheorem{theorem}{Theorem}
\newtheorem{lemma}{Lemma}
\newtheorem{corollary}{Corollary}
\theoremstyle{remark}
\newtheorem{remark}{Remark}
\newcommand{\T}{\mathbb T}
\newcommand{\Pol}{\operatorname{Pol}}
\newcommand{\Irr}{\operatorname{Irr}}
\newcommand{\id}{\mathrm{id}}
\newcommand{\Fix}{\operatorname{Fix}}

\title{All nonsymmetric convolution-power limits\\on a compact quantum group}
\author{Solution packet for arXiv:0802.1256}
\date{13 August 2026}

\begin{document}
\maketitle

\begin{abstract}
For an arbitrary state on an arbitrary compact quantum group, we classify the
weak-$*$ cluster set of its convolution powers.  It is a compact monothetic
group obtained by simultaneously rotating the peripheral eigenspaces of the
finite Peter--Weyl Fourier matrices.  This gives an exact criterion for
ordinary convergence, identifies the (possibly atypical) idempotent limit,
characterizes convergence to Haar, and yields point-norm convergence of the
associated convolution operators.  No symmetry, traciality, finite
dimensionality, or uniform spectral gap is required.
\end{abstract}

\section{The question and the answer}

Let $\mathbb G$ be a compact quantum group with reduced C$^*$-algebra $C(\mathbb G)$,
coproduct $\Delta$, and Haar state $h$.  For states $\omega,\rho$ write
$\omega\star\rho=(\omega\otimes\rho)\Delta$.  Franz and Skalski ask in
Section~4 of \cite{FS} for the limit behaviour of $\phi^{\star n}$, or of the
corresponding iterated convolution operators, when no symmetry assumption is
made on $\phi$.  The warning in their question is essential: an idempotent
state need not be the Haar state of a quantum subgroup.

Choose, for every $\alpha\in\Irr(\mathbb G)$, an irreducible unitary
corepresentation
\[
 U^\alpha=(u^\alpha_{ij})\in
 B(H_\alpha)\otimes\Pol(\mathbb G),
 \qquad
 \Delta(u^\alpha_{ij})=\sum_k u^\alpha_{ik}\otimes u^\alpha_{kj},
\]
and define the Fourier contraction
\[
 A_\alpha=(\id\otimes\phi)(U^\alpha)
       =\bigl(\phi(u^\alpha_{ij})\bigr)_{ij}\in B(H_\alpha).
\]
Let $\Lambda_\alpha=\sigma(A_\alpha)\cap\T$.  For
$\lambda\in\Lambda_\alpha$, let $P_{\alpha,\lambda}$ be the spectral
projection of $A_\alpha$ at $\lambda$.  These projections are orthogonal.
Put
\[
 I_\phi=\{(\alpha,\lambda):\lambda\in\Lambda_\alpha\},\qquad
 g_\phi=(\lambda)_{(\alpha,\lambda)\in I_\phi}\in\T^{I_\phi},
\]
and let
\[
 K_\phi=\overline{\{g_\phi^n:n\geq1\}}\subseteq\T^{I_\phi}.
\]
The closure is a compact monothetic group (and is the one-point group if
$I_\phi$ is empty).

\begin{theorem}[Complete Fourier--peripheral classification]\label{thm:main}
For each $z=(z_{\alpha,\lambda})\in K_\phi$ there is a unique state
$\phi_z$ on $C(\mathbb G)$ satisfying
\begin{equation}\label{eq:transform}
 (\id\otimes\phi_z)(U^\alpha)
 =B_\alpha(z):=
 \sum_{\lambda\in\Lambda_\alpha}
 z_{\alpha,\lambda}P_{\alpha,\lambda}
 \qquad(\alpha\in\Irr(\mathbb G)).
\end{equation}
Then:
\begin{enumerate}
\item The weak-$*$ cluster set of $(\phi^{\star n})_{n\geq1}$ is exactly
      $\{\phi_z:z\in K_\phi\}$.
\item The map $z\mapsto\phi_z$ is a topological group isomorphism from
      $K_\phi$ onto that cluster set equipped with convolution:
      $\phi_z\star\phi_w=\phi_{zw}$.  Its identity is
      $\phi_{\mathbf1}$, whose $\alpha$-Fourier matrix is
      $\sum_{\lambda\in\Lambda_\alpha}P_{\alpha,\lambda}$.
\item The Ces\`aro means $N^{-1}\sum_{n=1}^N\phi^{\star n}$ converge
      weak-$*$ to an idempotent state $\phi_{\mathrm C}$ characterized by
      $(\id\otimes\phi_{\mathrm C})(U^\alpha)=P_{\alpha,1}$.  It is absorbing
      for the cluster group:
      $\phi_z\star\phi_{\mathrm C}=\phi_{\mathrm C}
      =\phi_{\mathrm C}\star\phi_z$ for every $z\in K_\phi$.
\item The ordinary powers converge weak-$*$ if and only if
\begin{equation}\label{eq:criterion}
       \sigma(A_\alpha)\cap\T\subseteq\{1\}
       \quad\hbox{for every }\alpha\in\Irr(\mathbb G).
\end{equation}
      In that event $\phi^{\star n}\to\phi_{\mathbf1}=\phi_{\mathrm C}$ and
      $(\id\otimes\phi_{\mathbf1})(U^\alpha)=P_{\alpha,1}$.
\item One has $\phi^{\star n}\to h$ if and only if
\begin{equation}\label{eq:haar}
       \sigma(A_\alpha)\cap\T=\varnothing
       \quad\hbox{for every nontrivial }\alpha\in\Irr(\mathbb G).
\end{equation}
\end{enumerate}
For every convergence or cluster net above, the corresponding left or right
convolution operators converge point-norm on $C(\mathbb G)$.  In the reduced
Haar GNS realization they converge strongly on $L^2(\mathbb G,h)$.
\end{theorem}

Thus the all-compact answer is exact even when the limit is an atypical
idempotent.  Condition~\eqref{eq:criterion} is the aperiodicity condition;
condition~\eqref{eq:haar} adds ergodicity.

As a sanity check, take an ordinary compact group and $\phi=\delta_g$.
Then the cluster group consists of the point masses on
$\overline{\langle g\rangle}$, its identity is $\delta_e$, and
$\phi_{\mathrm C}$ is Haar measure on $\overline{\langle g\rangle}$.  This
also illustrates why the cluster-group identity and the Ces\`aro idempotent
must not be conflated.

\section{Finite contraction lemma}

\begin{lemma}\label{lem:contraction}
If $A$ is a contraction on a finite-dimensional Hilbert space, its
unit-circle eigenvalues are semisimple, their eigenspaces are mutually
orthogonal and reducing, and
\begin{equation}\label{eq:block}
 A^n-\sum_{\lambda\in\sigma(A)\cap\T}\lambda^nP_\lambda
 \longrightarrow0
\end{equation}
in norm.  Moreover $N^{-1}\sum_{n=1}^NA^n\to P_1$.
\end{lemma}

\begin{proof}
If $Ax=\lambda x$ and $|\lambda|=1$, then
$\langle(I-A^*A)x,x\rangle=0$.  Positivity gives $A^*Ax=x$, hence
$A^*x=\overline\lambda x$.  The eigenspace therefore reduces $A$; eigenspaces
for distinct unit-circle eigenvalues are orthogonal.  A nontrivial Jordan
block at a unit-circle eigenvalue would make $(A^n)$ unbounded, contrary to
$\|A^n\|\leq1$.  On the orthogonal complement of all these eigenspaces, the
spectrum lies strictly inside the unit disk, so the remaining powers tend to
zero.  This proves~\eqref{eq:block}.  Ces\`aro averaging retains only the
$\lambda=1$ summand.
\end{proof}

\section{Proof of the classification}

Slicing a unitary by a state is contractive, so $\|A_\alpha\|\leq1$.
The coproduct formula gives
\begin{equation}\label{eq:mult}
 (\id\otimes(\omega\star\rho))(U^\alpha)
 = (\id\otimes\omega)(U^\alpha)
   (\id\otimes\rho)(U^\alpha).
\end{equation}
In particular the Fourier matrix of $\phi^{\star n}$ is $A_\alpha^n$.
Lemma~\ref{lem:contraction} now gives, separately on every Peter--Weyl block,
\begin{equation}\label{eq:asymp}
 A_\alpha^n-
 \sum_{\lambda\in\Lambda_\alpha}
 \lambda^nP_{\alpha,\lambda}\longrightarrow0.
\end{equation}

The only point needing care is simultaneous passage over all irreducibles.
The closure of the positive powers of one element of a compact group is a
compact subgroup: its closure contains the identity (apply compactness to
pairs of sufficiently close powers), and hence contains inverses.  Thus
$K_\phi$ is a compact group.  Moreover every tail of $(g_\phi^n)$ is dense in
$K_\phi$, because its closure is $g_\phi^N K_\phi=K_\phi$.  Thus every point
of $K_\phi$ is realized by a subnet with $n\to\infty$.

Fix $z\in K_\phi$ and choose a net $(n_i)$ with $g_\phi^{n_i}\to z$.
By~\eqref{eq:asymp}, the Fourier matrices of $\phi^{\star n_i}$ converge to
$B_\alpha(z)$ for every $\alpha$.  The state space is weak-$*$ compact, so a
subnet has a weak-$*$ limit.  Every such limit has the Fourier matrices
in~\eqref{eq:transform}; since the coefficients $u^\alpha_{ij}$ span the
norm-dense Hopf $*$-algebra $\Pol(\mathbb G)$, these matrices determine the
functional uniquely.  Hence the original net converges to a unique state
$\phi_z$.  The same density argument shows that $z\mapsto\phi_z$ is continuous
and injective.

Conversely, from any weak-$*$ cluster net of $(\phi^{\star n})$, compactness of
$K_\phi$ supplies a further subnet on which $g_\phi^n$ converges to some $z$;
the preceding paragraph identifies the cluster state with $\phi_z$.  This
proves part~1.  Orthogonality of the peripheral spectral projections and
\eqref{eq:mult} give
\[
 B_\alpha(z)B_\alpha(w)=B_\alpha(zw),
\]
so Fourier uniqueness yields $\phi_z\star\phi_w=\phi_{zw}$.  Compactness then
gives the topological group assertion in part~2.

By Lemma~\ref{lem:contraction}, the Fourier matrices of the Ces\`aro means
converge to $P_{\alpha,1}$.  Norm density and uniform boundedness of the states
give weak-$*$ convergence to a state $\phi_{\mathrm C}$ with those Fourier
matrices.  Since $P_{\alpha,1}^2=P_{\alpha,1}$, Fourier uniqueness proves
idempotence.  The coordinate belonging to the phase $\lambda=1$ is identically
one throughout $K_\phi$, so
$B_\alpha(z)P_{\alpha,1}=P_{\alpha,1}=P_{\alpha,1}B_\alpha(z)$; this proves the
two absorption identities.  Notice that the cluster-group identity
$\phi_{\mathbf1}$ retains \emph{all} peripheral eigenspaces, whereas Ces\`aro
averaging retains only the phase-one eigenspace.  They coincide precisely when
condition~\eqref{eq:criterion} holds.

The scalar sequence $\lambda^n$, $\lambda\in\T$, converges only when
$\lambda=1$.  Equations~\eqref{eq:asymp} and density therefore prove
part~4.  Finally, the Fourier matrix of the Haar state is $1$ on the trivial
corepresentation and $0$ on every nontrivial irreducible.  This proves
part~5.

For the operator claim, on each coefficient space convolution is multiplication
by the relevant Fourier matrix.  Hence the asserted operator convergence is
norm convergence on $\Pol(\mathbb G)$.  Convolution by a state is contractive;
approximating an arbitrary element by coefficient polynomials proves
point-norm convergence on $C(\mathbb G)$.  The Schwarz inequality and Haar
preservation give contractions on the Haar GNS $L^2$ space, and density of the
coefficient algebra proves strong convergence there.

\section{A support-free periodicity test}

The peripheral obstruction admits an intrinsic equality-case description.
Let $(\pi_\phi,H_\phi,\Omega_\phi)$ be the GNS representation of $\phi$ and
put
\[
 V_\alpha=(\id\otimes\pi_\phi)(U^\alpha),\qquad
 J_\alpha\xi=\xi\otimes\Omega_\phi.
\]
Then $A_\alpha=J_\alpha^*V_\alpha J_\alpha$.

\begin{corollary}[Deterministic-vector criterion]\label{cor:gns}
For $\lambda\in\T$ and $0\ne\xi\in H_\alpha$,
\[
 A_\alpha\xi=\lambda\xi
 \quad\Longleftrightarrow\quad
 V_\alpha(\xi\otimes\Omega_\phi)
 =\lambda\xi\otimes\Omega_\phi.
\]
Consequently $\phi^{\star n}$ converges precisely when every deterministic
corepresentation vector has phase $1$, and it converges to Haar precisely
when no nontrivial irreducible corepresentation has any deterministic vector.
\end{corollary}

\begin{proof}
The reverse implication follows by compression.  For the forward implication,
\[
 \|J_\alpha^*V_\alpha J_\alpha\xi\|=\|\xi\|
 =\|V_\alpha J_\alpha\xi\|.
\]
Equality for the orthogonal projection $J_\alpha J_\alpha^*$ forces
$V_\alpha J_\alpha\xi$ to lie in the range of $J_\alpha$.  Its compression is
$\lambda\xi$, so it equals $\lambda J_\alpha\xi$.
\end{proof}

For an ordinary compact group, this says that a unitary representation vector
is multiplied by one fixed phase almost everywhere on the support of the
driving measure---the familiar periodic-coset obstruction.  In the quantum
case it remains meaningful without forcing the limiting idempotent to arise
from a subgroup.

\begin{remark}
McCarthy's finite-dimensional theorem \cite{McC} translates irreducibility and
periodicity into quasi-subgroups and cyclic cosets.  Theorem~\ref{thm:main}
instead works for every compact quantum group and records all cluster states,
including infinitely many phases and non-Haar idempotent limits.  Related
fixed-point and Ces\`aro theory for locally compact quantum groups appears in
\cite{NSSS}.
\end{remark}

\begin{thebibliography}{9}
\bibitem{FS}
U. Franz and A. Skalski,
\emph{On ergodic properties of convolution operators associated with compact
quantum groups}, Colloq. Math. 113 (2008), 13--23; arXiv:0802.1256.

\bibitem{McC}
J. P. McCarthy,
\emph{The ergodic theorem for random walks on finite quantum groups},
Commun. Algebra 49 (2021), 3850--3871; arXiv:2004.01234.

\bibitem{NSSS}
M. Neufang, P. Salmi, A. Skalski, and N. Spronk,
\emph{Fixed points and limits of convolution powers of contractive quantum
measures}, Indiana Univ. Math. J. 70 (2021), 1971--2009;
arXiv:1907.07337.
\end{thebibliography}

\end{document}
